Stealthy false data injection attacks (FDIAs) shift the estimated state of a power system without increasing the state-estimation residual, so the conventional $\chi^2$ bad-data detector cannot see them. This paper proposes DT-GCN, a framework that combines a physics-based digital twin (DT) with a graph convolutional network to detect and localize such attacks. The twin runs AC power flow on very-short-term injection forecasts and schedules, which gives an operating-state reference that does not use the possibly compromised telemetry. Bus-level twin-residual features are processed on the grid topology by a dual-head GCN that outputs an attack probability for the snapshot and for each bus. We evaluate the method on the IEEE 14- and 118-bus systems using AC power flow, nonlinear weighted-least-squares estimation, 15-min load and photovoltaic profiles, and AC-stealthy attacks. The $\chi^2$ test detects only %%BDDST118%%\% of these attacks. On the 118-bus system, DT-GCN reaches %%F118%%\% detection F1 at %%FPR118%%\% false-positive rate, against at most %%GCNRAW118%%\% for SVM, random forest, MLP and GCN detectors trained on raw telemetry. It identifies the exact set of compromised buses in %%LOCEX118%%\% of attacks, with %%LOCF118%%\% bus-level F1. Inference takes under 0.1~ms per snapshot, and the code is released for reproducibility.
